% SEGMENT OLP-0363-S01
% Part: lambda-calculus
% Chapter: syntax
% Section: de-bruijn

\documentclass[../../../include/open-logic-section]{subfiles}

\begin{document}

\olfileid{lam}{syn}{deb}

\olsection{டெ புரூய்ன் சுட்டெண்}

$\alpha$-சமானம் மிகவும் இயல்பானது; ஏனெனில் $\alpha$-சமானமான சொற்கள்
``ஒரே பொருளைக் கொண்டுள்ளன.'' ஒன்றிலிருந்து மற்றொன்றுக்கு $\alpha$-மாற்றம்
செய்யக்கூடிய சொற்களை வேறுபடுத்தாத லாம்டா சொல் தொடரமைப்பை உண்மையில் வழங்க
முடியும். அவற்றுள் மிக அறியப்பட்டது மாறிகளை அவற்றின் \emph{டெ புரூய்ன்
சுட்டெண்களால்} மாற்றுகிறது.

$\lambd[x][M]$ என்று எழுதும்போது, $x$ ஆனது சார்பின் அளபுரு என்பதை
வெளிப்படையாகக் கூறுகிறோம்; எனவே இந்த அளபுருவைச் சுட்டுவதற்கு $M$ இல்
$x$ ஐப் பயன்படுத்தலாம். ஆனால் டெ புரூய்ன் சுட்டெண் முறையில் அளபுருக்கள்
பெயரற்றவை; சார்பின் உடலில் அவற்றைச் சுட்டுவது, அந்த இடத்திற்கும் அதைக் கட்டும்
அருவமாக்கத்திற்கும் இடையிலுள்ள அருவமாக்க நிலைகளின் எண்ணிக்கையால்
குறிக்கப்படுகிறது. எடுத்துக்காட்டாக, $\lambd[x][\lambd[y][y x]]$ ஐக்
கருதுக: வெளிப்புற அருவமாக்கம் மாறி~$x$ ஐக் கட்டுகிறது; உட்புற அருவமாக்கம்
மாறி~$y$ ஐக் கட்டுகிறது; உட்சொல் $y x$ உட்புற அருவமாக்கத்தின் வீச்சில் உள்ளது.
$y$ க்கும் அதைக் கட்டும் $\lambd[y]$ அருவத்திற்கும் இடையில் அருவமாக்கம்
எதுவுமில்லை; $x$ க்கும் அதைக் கட்டும் $\lambd[x]$ அருவத்திற்கும் இடையில் ஓர்
அருவமாக்கம் உள்ளது. ஆகவே $y x$ க்குப் பதிலாக $0\, 1$ என்றும், முழுச்
சொல்லுக்குப் பதிலாக $\lambd[][\lambd[][01]]$ என்றும் எழுதுகிறோம்.
% SOURCE CORRECTION TA-LAM-020: Tamil resolves the frozen malformed binder sentences and ordinary prose inflection errors while retaining the stated zero-based indexing convention.

\begin{defn}
  டெ புரூய்ன் சொற்கள் பின்வருமாறு தொகுத்தறிதல் முறையில் வரையறுக்கப்படுகின்றன:
  \begin{enumerate}
  \item $n$; இங்கு $n$ ஏதேனும் இயல் எண்.
  \item $PQ$; இங்கு $P$, $Q$ இரண்டும் டெ புரூய்ன் சொற்கள்.
  \item $\lambd[][N]$; இங்கு $N$ ஒரு டெ புரூய்ன் சொல்.
  \end{enumerate}
\end{defn}

சாதாரண லாம்டா சொற்களிலிருந்து டெ புரூய்ன் சுட்டெண் கொண்ட சொற்களுக்கு ஒரு
முறைப்படுத்தப்பட்ட மாற்றம் பின்வருமாறு:
\begin{defn}
  \begin{align*}
    F_\Gamma(x) &= \Gamma(x) \\
    F_\Gamma(PQ) &= F_\Gamma(P)F_\Gamma(Q) \\
    F_\Gamma(\lambd[x][N]) &= \lambd[][F_{x,\Gamma}(N)]
  \end{align*}
  இங்கு $\Gamma$ என்பது பூச்சியத்திலிருந்து சுட்டெண்கள் இடப்பட்ட மாறிகளின்
  பட்டியல்; $\Gamma(x)$ என்பது $\Gamma$ இல் மாறி $x$ இன் இடத்தைக் குறிக்கிறது.
  எடுத்துக்காட்டாக, $\Gamma$ என்பது $x,y,z$ எனில், $\Gamma(x)$ என்பது $0$;
  $\Gamma(z)$ என்பது $2$.

  $x,\Gamma$ என்பது $\Gamma$ இன் தலைப்பகுதியில் $x$ ஐச் சேர்ப்பதால் கிடைக்கும்
  பட்டியலைக் குறிக்கிறது; எடுத்துக்காட்டாக, முந்தைய எடுத்துக்காட்டின் $\Gamma$
  ஐத் தொடர்ந்தால், $w,\Gamma$ என்பது $w,x,y,z$.
\end{defn}

ஒரு டெ புரூய்ன் சொல்லிலிருந்து தரநிலை லாம்டா சொல்லை மீட்பது பின்வருமாறு
செய்யப்படுகிறது:

\begin{defn}
  \begin{align*}
    G_\Gamma(n) &= \Gamma[n] \\
    G_\Gamma(PQ) &= G_\Gamma(P) G_\Gamma(Q) \\
    G_\Gamma(\lambd[][N]) &= \lambd[x][G_{x,\Gamma}(N)]
  \end{align*}
  இங்கும் $\Gamma$ என்பது பூச்சியத்திலிருந்து சுட்டெண்கள் இடப்பட்ட மாறிகளின்
  பட்டியல்; $\Gamma[n]$ என்பது $n$ ஆம் இடத்திலுள்ள மாறியைக் குறிக்கிறது.
  எடுத்துக்காட்டாக, $\Gamma$ என்பது $x,y,z$ எனில், $\Gamma[1]$ என்பது $y$.

  கடைசிச் சமன்பாட்டில் உள்ள மாறி $x$, $\Gamma$ இல் இல்லாத ஏதேனும் ஒரு
  மாறியாகத் தேர்ந்தெடுக்கப்படுகிறது.
\end{defn}

சில முடிவுகளை நிறுவாமல் இங்குத் தருகிறோம்:

\begin{prop}
  $M \aconvone M'$ ஆகவும், $\Gamma$ ஆனது $\FV{M}$ ஐக் கொண்ட ஏதேனும் ஒரு
  பட்டியலாகவும் இருந்தால், $F_\Gamma(M) \eqs F_\Gamma(M')$.
\end{prop}

\end{document}
